\section{Cut clearance of belt and descendant birth templates}
\label{sec:area_law_descendant_cut_clearance}

The belt and replacement templates of \cite[Section~11]{OpenAI2026AreaLaw}
remain far from the cut endpoints $Z$ compared with their size. Write
$r=2^k$ and $t_k$ for the layer and fine sides of
Definition~\ref{def:area_law_fine_belt_scales}, $C_0\ge2$ for the layer
width, and $d(S)$ for the removal half-side of
Definition~\ref{def:area_law_repair_square}. A belt tile has side scale
$u=t_k$, piece size $s_0=u$ and scanner scale $n=\lceil C_Bu\rceil$; a
replacement whose removal square has half-side $u$ has $s_0=2u$ and
$n=\lceil C_Bu\rceil$. The constants $D_0$ and $C_B$ are arbitrary, and the
lower layer index $k_0$ is increased after them.

% Source: 10-geometry.tex, lines 742--750, with the scanner scales of
% lines 732--737 and the order of constants, lines 756--774.
% Revision: adc7f1241b42e322a6451854ab7e4b4c146bf78a.

\begin{theorem}[Scale inequalities of late layers]
    \label{thm:area_law_descendant_clearance_scales}
    \lean{TNLean.PEPS.AreaLaw.Geometry.two_pow_fineScaleIndex_le_rpow,
        TNLean.PEPS.AreaLaw.Geometry.eventually_descendant_cutClearance_scales}
    \leanok
    \uses{def:area_law_fine_belt_scales,thm:area_law_exponent_gaps,
        thm:area_law_repair_margins}
    Always $t_k\le r^\zeta$. For all real $D_0$ and $C_B$, every
    sufficiently large $k$ satisfies
    \begin{align}
        2t_k^{\rho_0}              & \le r, &
        4D_0\cdot2t_k              & <r,    &
        2\lceil C_Bt_k\rceil^\beta & <r.
    \end{align}
\end{theorem}
\begin{proof}
    \leanok
    \uses{def:area_law_fine_belt_scales,thm:area_law_exponent_gaps,
        thm:area_law_repair_margins}
    The first claim is $\lfloor\zeta k\rfloor\le\zeta k$. For a fixed $M$
    and an exponent $a<1$, we have $Mr^a=Mr^{-(1-a)}r<r$ once
    $r^{-(1-a)}<1/(\max\{M,0\}+1)$. Apply this with $a=\zeta\rho_0$,
    $a=\zeta$ and $a=\zeta\beta$, all less than one because $\zeta$,
    $\rho_0$ and $\beta$ lie in $[0,1)$. For the third inequality put
    $B=\max\{C_B,0\}+1$; since $t_k\ge1$,
    $\lceil C_Bt_k\rceil\le Bt_k$.
\end{proof}

\begin{theorem}[Template clearances from endpoint distance]
    \label{thm:area_law_template_clearance_of_endpoint_distance}
    \lean{TNLean.PEPS.AreaLaw.Geometry.Template.cutClearance_of_endpoint_dist_ge}
    \leanok
    \uses{def:area_law_template,thm:area_law_descendant_clearance_scales}
    Let $T$ be a template with parameters $n$ and $s_0$ whose points all
    have endpoint distance at least $(C_0-1)r$. Suppose that
    $0\le u\le t_k$, $s_0\le2u$ and $n\le\lceil C_Bu\rceil$, and that the
    second and third inequalities of
    Theorem~\ref{thm:area_law_descendant_clearance_scales} hold at $k$.
    Then
    \begin{align}
        \operatorname{dist}_\infty(T,Z) & >4D_0s_0,  &
        \operatorname{dist}_\infty(T,Z) & >2n^\beta.
    \end{align}
\end{theorem}
\begin{proof}
    \leanok
    Since $C_0\ge2$, the endpoint distance is at least $r$. If $D_0\ge0$,
    then $4D_0s_0\le4D_0\cdot2t_k<r$; otherwise $4D_0s_0\le0<r$. Also
    $n\le\lceil C_Bt_k\rceil$, because the ceiling is monotone when
    $C_B\ge0$ and $\lceil C_Bu\rceil=0$ when $C_B<0$. Hence
    $2n^\beta\le2\lceil C_Bt_k\rceil^\beta<r$.
\end{proof}

\begin{theorem}[Initial marks are far from the cut]
    \label{thm:area_law_initial_mark_endpoint_distance}
    \lean{TNLean.PEPS.AreaLaw.Geometry.repairStage_zero_endpoint_dist_ge}
    \leanok
    \uses{thm:area_law_initial_repair_mark_facts,
        thm:area_law_belt_marks_closed_cell,
        thm:area_law_fine_cell_containment,
        thm:area_law_dyadic_layer_distance}
    Every initial mark $v$ of scale $S_v$ is a mark of a belt cell of some
    layer $k\ge k_0$ with $S_v=t_k$, and
    \begin{align}
        \operatorname{dist}_\infty(v,Z) & \ge C_0r.
    \end{align}
\end{theorem}
\begin{proof}
    \leanok
    \uses{thm:area_law_initial_repair_mark_facts,
        thm:area_law_belt_marks_closed_cell,
        thm:area_law_fine_cell_containment,
        thm:area_law_dyadic_layer_distance}
    The initial mark facts supply the layer $k$ and the belt cell whose
    mark is $v$ and whose fine side is $S_v$. The mark lies in the closed
    belt cell, which lies in $\overline{D_k}$. The layer distance estimate
    gives the bound.
\end{proof}

\begin{theorem}[Descendant birth regions are far from the cut]
    \label{thm:area_law_descendant_birth_endpoint_distance}
    \lean{TNLean.PEPS.AreaLaw.Geometry.repairStage_birth_endpoint_dist_ge}
    \leanok
    \uses{thm:area_law_initial_mark_endpoint_distance,
        def:area_law_uniform_repair_conclusion,
        thm:area_law_repair_descendant_scales}
    Assume the descendant bound of the uniform repair induction. Let a
    replacement of generation $j+1$ have parent mark $m$ of scale $S_m$,
    and let its initial ancestor have scale $t_k$, with $k\ge k_0$. Then
    $d(S_m)\le t_k$, and every point $y$ of its birth region satisfies
    \begin{align}
        \operatorname{dist}_\infty(y,Z) & \ge C_0r-2d(t_k).
    \end{align}
\end{theorem}
\begin{proof}
    \leanok
    \uses{thm:area_law_initial_mark_endpoint_distance,
        def:area_law_uniform_repair_conclusion,
        thm:area_law_repair_descendant_scales,thm:area_law_repair_margins}
    The scale $S_m$ is the $j$-th iterate of $d$ at $t_k$. For $S\ge1$ we
    have $1\le d(S)\le S^{\rho_0}\le S$, so $d(S_m)\le t_k$. The birth
    region lies within $2d(t_k)$ of the ancestor mark, and the triangle
    inequality with
    Theorem~\ref{thm:area_law_initial_mark_endpoint_distance} gives the
    bound. No layer of the birth region is used, so this also covers
    replacements extending into the dummy region.
\end{proof}

\begin{theorem}[Cut clearance of belt birth templates]
    \label{thm:area_law_belt_birth_template_clearance}
    \lean{TNLean.PEPS.AreaLaw.Geometry.initialBelt_birthTemplate_cutClearance}
    \leanok
    \uses{thm:area_law_template_clearance_of_endpoint_distance,
        thm:area_law_fan_run_cover,thm:area_law_fine_cell_containment,
        thm:area_law_dyadic_layer_distance,
        def:area_law_initial_birth_regions}
    For all real $D_0$ and $C_B$ there is a threshold such that, whenever
    $k_0$ exceeds it, every belt tile of a layer $k\ge k_0$ has birth
    region at endpoint distance at least $C_0r$, and every template $T$
    whose points lie in that birth region, with $s_0\le2t_k$ and
    $n\le\lceil C_Bt_k\rceil$, satisfies
    \begin{align}
        \operatorname{dist}_\infty(T,Z) & >4D_0s_0,  &
        \operatorname{dist}_\infty(T,Z) & >2n^\beta.
    \end{align}
    This covers the belt choice $s_0=t_k$ and $n=\lceil C_Bt_k\rceil$.
\end{theorem}
\begin{proof}
    \leanok
    \uses{thm:area_law_fan_run_cover,thm:area_law_fine_cell_containment,
        thm:area_law_dyadic_layer_distance,
        thm:area_law_descendant_clearance_scales,
        thm:area_law_template_clearance_of_endpoint_distance}
    The birth region of a belt tile is one of the runs covering its closed
    belt cell, which lies in $\overline{D_k}$; the layer distance estimate
    gives $C_0r\ge(C_0-1)r$. Take the threshold of
    Theorem~\ref{thm:area_law_descendant_clearance_scales} and apply
    Theorem~\ref{thm:area_law_template_clearance_of_endpoint_distance}
    with $u=t_k$.
\end{proof}

\begin{theorem}[Cut clearance of descendant birth templates]
    \label{thm:area_law_descendant_birth_template_clearance}
    \lean{TNLean.PEPS.AreaLaw.Geometry.descendant_birthTemplate_cutClearance}
    \leanok
    \uses{thm:area_law_uniform_repair_induction,
        thm:area_law_descendant_birth_endpoint_distance,
        thm:area_law_descendant_clearance_scales,
        thm:area_law_template_clearance_of_endpoint_distance}
    There is a threshold for the stopping scale $S_*$ such that the
    following holds for every $S_*$ beyond it, every $C_0\ge2$ and all
    real $D_0$ and $C_B$. For every sufficiently large $k_0$, each
    replacement of each generation, with removal half-side $u$, has an
    initial ancestor in a layer $k\ge k_0$ with $u\le t_k$. Its birth
    region has endpoint distance at least $(C_0-1)r$, including when it
    extends into the dummy region. Every template $T$ whose points lie in
    that birth region, with $s_0\le2u$ and $n\le\lceil C_Bu\rceil$,
    satisfies
    \begin{align}
        \operatorname{dist}_\infty(T,Z) & >4D_0s_0,  &
        \operatorname{dist}_\infty(T,Z) & >2n^\beta.
    \end{align}
    This covers the replacement choice $s_0=2u$ and
    $n=\lceil C_Bu\rceil$.
\end{theorem}
\begin{proof}
    \leanok
    \uses{thm:area_law_uniform_repair_induction,
        thm:area_law_descendant_clearance_scales,
        thm:area_law_descendant_birth_endpoint_distance,
        thm:area_law_repair_margins,
        thm:area_law_template_clearance_of_endpoint_distance}
    Take the thresholds of the uniform repair induction for $S_*$ and of
    Theorem~\ref{thm:area_law_descendant_clearance_scales} for $k_0$. By
    Theorem~\ref{thm:area_law_descendant_birth_endpoint_distance} and
    $2d(t_k)\le2t_k^{\rho_0}\le r$, the birth region has endpoint distance
    at least $C_0r-r=(C_0-1)r$. Then apply
    Theorem~\ref{thm:area_law_template_clearance_of_endpoint_distance}.
\end{proof}
